% !TeX root = ../../Book3.tex
% 纲要标题：### 5.4 Noether 正规化
% 纲要位置：工作记录/计划纲要.md，第 2054 行。
% 纲要细目（写作范围）：
% * 域上有限生成代数的整子多项式环；与分式域内整闭包的区别
% * 统一的加权证明与无限域上的线性代换；唯一最高幂与有限域的取值障碍
% * 超越次数、整环情形及有限态射的纤维；分式域识别与有限维 Artin 纤维
% ---

\section{Noether 正规化}\label{b3:sec:0504}

有限型代数通常由带关系的生成元给出。Noether 正规化要从中找出一组没有代数关系的元素，使原代数在所选多项式子环上成为有限生成模。这一构造既为维数提供坐标，也为下一章的正规化有限性提供一个可控制的底环。

\subsection{有限生成代数的整子多项式环}\label{b3:subsec:0504:01}

\begin{theorem}[Noether 正规化]\label{b3:thm:noether-normalization}
设 \(k\) 为任意域，\(A\ne0\) 为交换含幺有限型 \(k\) 代数。存在代数独立元素 \(y_1,\ldots,y_d\in A\)，使 \(A\) 是多项式子环
\[
R=k[y_1,\ldots,y_d]\subseteq A
\]
上的有限 \(R\) 代数；这里允许 \(d=0\)，此时底环为 \(k\)。这称为\term[noether-zhengguihua]{Noether 正规化定理}[Noether normalization theorem]。
\end{theorem}

本定理不要求 \(A\) 是整环或约化环。例如
\(k[x,\varepsilon]/(\varepsilon^2)\) 作为 \(k[x]\) 模由 \(1,\varepsilon\) 生成，\(x\) 本身仍然代数独立。对整环而言，选取多项式子环与在分式域内取整闭包是两个不同的构造，\cref{b3:tab:two-normalizations}列出了它们的区别。

\begin{table}[htbp]
\centering
\caption[Noether 正规化与环的正规化]{Noether 正规化与环的正规化。设 \(k\) 为域，\(A\) 为交换含幺有限型 \(k\) 整环，\(K=\operatorname{Frac}(A)\)，\(\widetilde A\) 为 \(A\) 在 \(K\) 中的整闭包。}
\label{b3:tab:two-normalizations}
\begin{tabularx}{\textwidth}{@{}>{\raggedright\arraybackslash}X>{\raggedright\arraybackslash}X@{}}
\toprule
Noether 正规化 & 环的正规化\\
\midrule
选出多项式子环 \(R=k[y_1,\ldots,y_d]\subseteq A\) & 添入 \(K\) 中所有在 \(A\) 上整的元素，得到 \(\widetilde A\)\\
\(A\) 是有限 \(R\) 代数 & \(\widetilde A\) 在 \(A\) 上整\\
\(R\) 是正规整环，\(A\) 未必正规 & \(\widetilde A\) 是分式域为 \(K\) 的正规整环\\
保留 \(A\) 本身，选择其子环 & \(A\subseteq\widetilde A\)，且两者相等当且仅当 \(A\) 正规\\
\bottomrule
\end{tabularx}
\end{table}

\subsection{统一的加权证明与无限域上的线性代换}\label{b3:subsec:0504:02}

\begin{proof}[\cref{b3:thm:noether-normalization}的证明]
写 \(A=k[a_1,\ldots,a_n]\)，对生成元个数 \(n\) 归纳。\(n=0\) 时 \(A=k\)。若 \(a_1,\ldots,a_n\) 代数独立，则已完成构造。否则存在非零多项式 \(f\in k[X_1,\ldots,X_n]\)，使 \(f(a_1,\ldots,a_n)=0\)。若能改选前 \(n-1\) 个生成元，使这个关系成为关于 \(a_n\) 的首一方程，\(a_n\) 就整于前面新生成元组成的子代数，而 \(A\) 在该子代数上有限。这样便能把归纳应用于一个少用一个生成元的子代数。下面的代换正是为使最高次项的系数成为 \(k\) 中的非零常数。

选整数 \(N\) 大于 \(f\) 中出现的全部单个变量指数，对 \(i<n\) 令 \(m_i=N^i\)，并作代换
\[
X_i=Y_i+Y_n^{m_i}\quad(i<n),\qquad X_n=Y_n.
\]
原来指数为 \(\alpha=(\alpha_1,\ldots,\alpha_n)\) 的单项式，在新的 \(Y_n\) 中产生最高次数
\[
w(\alpha)=\alpha_n+\alpha_1N+\cdots+\alpha_{n-1}N^{n-1}.
\]
由于每个指数小于 \(N\)，不同的 \(\alpha\) 有不同的 \(N\) 进制表示，从而这些次数互异。取其中最大的一个，其最高次 \(Y_n\) 项的系数恰为原单项式的非零系数；其他项不能抵消它。例如，对二变量多项式
\[
X_1^2+X_1X_2^2+X_2,
\]
取 \(N=3\)，三项的权重分别为 \(6,5,1\)。代入 \(X_1=Y_1+Y_2^3\)、\(X_2=Y_2\) 后，这三项产生的最高 \(Y_2\) 次数分别就是 \(6,5,1\)，所以 \(Y_2^6\) 只有一个来源，不会被抵消。

回到所取的关系，将 \(f\) 除以变换后最高次项的非零系数，并仍记为 \(f\)。令
\[
g(Y_1,\ldots,Y_n)
=f(Y_1+Y_n^{m_1},\ldots,Y_{n-1}+Y_n^{m_{n-1}},Y_n).
\]
于是 \(g\) 作为 \(Y_n\) 的多项式首一，系数在 \(k[Y_1,\ldots,Y_{n-1}]\) 中。令 \(b_i=a_i-a_n^{m_i}\)，实际代入给出
\[
g(b_1,\ldots,b_{n-1},a_n)=f(a_1,\ldots,a_n)=0.
\]
这是一条 \(C=k[b_1,\ldots,b_{n-1}]\) 系数首一方程，所以 \(a_n\) 在 \(C\) 上整；又 \(a_i=b_i+a_n^{m_i}\)，故 \(A=C[a_n]\) 是有限 \(C\) 代数。\(C\) 包含 \(k\)，因而非零，且至多由 \(n-1\) 个元素作为 \(k\) 代数生成；所减少的是选定的代数生成元个数。归纳假设给出一个多项式子环 \(R=k[y_1,\ldots,y_d]\)，使 \(C\) 为有限 \(R\) 代数。将两层有限组模生成元相乘便知 \(A\) 也是有限 \(R\) 代数，证明完成。
\end{proof}

加权代换对有限域和无限域同样有效。无限域上还有次数通常较低的做法：设 \(f\) 的最高齐次部分为 \(f_m\)。非零多项式在无限域上不可能处处为零，因此可取
\[
f_m(c_1,\ldots,c_{n-1},1)\ne0.
\]
这里代入最后变量为一不会把非零齐次式变成零多项式：其余各变量的指数确定后，最后一个指数由总次数唯一确定。作代换 \(X_i=Y_i+c_iY_n\) 后，\(Y_n^m\) 的系数为上述非零值；除以这一系数，就使变换后的关系对 \(Y_n\) 首一。所用的非零多项式取值事实可对变量个数归纳：固定前面的变量使某个非零系数不消失，再避开最后一个变量的有限个根。

\ProseParagraphBreak
有限域上这一步不能照搬。例如
\(X^qY-XY^q\) 是 \(\mathbb F_q[X,Y]\) 中的非零齐次多项式，却在所有 \((a,b)\in\mathbb F_q^2\) 处取零，因为 \(a^q=a\)、\(b^q=b\)。因此不能保证选到一个域内取值使最高齐次部分非零。上面的加权证明只比较多项式的指数，不需要寻找非零取值，因而同样适用于有限域。

\subsection{超越次数与整环情形}\label{b3:subsec:0504:03}

设 \(k\) 为域，\(A\) 为交换含幺有限型 \(k\) 整环。Noether 正规化得到的多项式子环 \(R=k[y_1,\ldots,y_d]\) 满足
\[
k(y_1,\ldots,y_d)\subseteq\operatorname{Frac}(A)
\]
为有限域扩张。事实上，记 \(K=\operatorname{Frac}(R)\)、\(S=R\setminus\{0\}\)。局部化 \(S^{-1}A\cong A\otimes_RK\) 是有限维 \(K\) 代数，因而在 \(K\) 上整；它又是整环，由\cref{b3:lem:integral-field-criterion}可知它是域。由于 \(R\subseteq A\) 都是整环，\(S^{-1}A\) 自然包含于 \(\operatorname{Frac}(A)\)，且包含 \(A\)。既然它是域，\(A\) 中每个非零元素的逆也在其中，故反过来 \(\operatorname{Frac}(A)\subseteq S^{-1}A\)，两者相等。所以 \(d\) 等于此分式域在 \(k\) 上的超越次数：\(y_i\) 代数独立，整个扩域在它们生成的有理函数域上代数。这里使用第一卷定义14.3.1的超越基判据及定理14.3.6的超越次数不变性；本卷后面的维数理论将把这个数与素理想链长度联系起来。

\begin{example}\label{b3:exm:noether-normalization-curves}
设 \(k\) 为域，\(A=k[x,y]/(xy-1)\)，并以 \(x,y\) 仍记其剩余类。按 \(x\) 坐标将双曲线 \(xy=1\) 投影到 \(x\) 轴，对应的环包含为 \(k[x]\subseteq k[x,x^{-1}]=A\)。\(A\) 不是有限 \(k[x]\) 代数，因为 \(y=x^{-1}\) 不是在 \(k[x]\) 上的整元素。改用投影坐标 \(u=x+y\) 则不同：
\[
x^2-ux+1=0,\qquad y=u-x.
\]
因此 \(A=k[u][x]\) 由 \(1,x\) 作为 \(k[u]\) 模生成。给定一个 \(u\) 值，\(x\) 必须满足上面的首一二次方程，\(y\) 再由 \(y=u-x\) 确定；在任意扩域中，这至多给出两个点。换投影方向后，整性用这一条首一关系控制了剩余坐标，使坐标环在 \(k[u]\) 上有限。\(u\) 代数独立，因为在同构 \(A=k[x,x^{-1}]\) 下，若次数 \(r\) 的非零多项式在 \(x+x^{-1}\) 处取零，其最高 Laurent 次数 \(x^r\) 的系数却非零，矛盾。对尖点代数 \(k[t^2,t^3]\)，直接取 \(u=t^2\) 即得一个由 \(1,t^3\) 作为 \(k[u]\) 模生成的有限 \(k[u]\) 代数。
\end{example}

\subsection{有限态射与纤维}\label{b3:subsec:0504:04}

\begin{definition}\label{b3:def:finite-affine-map}
设 \(R,A\) 为交换含幺环，\(\varphi:R\to A\) 为保单位环同态。若 \(A\) 经 \(\varphi\) 成为有限生成 \(R\) 模，则称诱导的素谱映射 \(\operatorname{Spec}A\to\operatorname{Spec}R\) 为\term[youxian-taishe]{有限态射}[finite morphism]。
\end{definition}

这一条件要求选定有限组元素 \(h_1,\ldots,h_r\in A\)，使 \(A\) 中的每个元素都能写成 \(\sum_i\varphi(c_i)h_i\)，其中 \(c_i\in R\)。因此它控制整个坐标环的模结构。上例中双曲线到 \(x\) 轴的投影对应 \(k[x]\subset k[x,x^{-1}]\)，其素谱映射是单射，每个纤维至多一个点，坐标环却不是有限 \(k[x]\) 模。所以只数点集纤维不能判定有限态射。

\ProseParagraphBreak
Noether 正规化中 \(R=k[y_1,\ldots,y_d]\subseteq A\) 的包含诱导有限态射
\[
\operatorname{Spec}A\longrightarrow\operatorname{Spec}R.
\]
由于这个环同态单射且 \(A\) 在 \(R\) 上整，\cref{b3:thm:lying-over}保证上述态射满射。一般的有限态射也有有限的点集纤维：对定义中的任意 \(\mathfrak p\in\operatorname{Spec}R\)，纤维对应于
\[
\operatorname{Spec}\bigl(A\otimes_R\kappa(\mathfrak p)\bigr).
\]
该纤维代数在剩余域上有限维，所以理想降链中每次严格包含都会使向量空间维数下降，降链必在有限步后稳定；因此它是 Artin 环。若它为零代数，纤维为空；否则由\cref{b3:lem:artin-primes-radical}，其素理想全为极大理想，且只有有限多个。这就证明了点集纤维的有限性，所用的纤维对应正是\cref{b3:prop:fibre-spectrum}。有限态射仍可能把不同点合在一起，也可能有非约化纤维；有限性没有消除这些现象。
